%------------------------------------------------------------------------------%
\begin{question}
  Dada a função
  \[
    f(x, y) = 6x - 3y + 9
  \]

  a) Determine e esboce a curva de nível \(f(x, y) = 6\)

  b) Calcule e esboce o gradiente de $f$ nos pontos
  \((0, 1)\) e
  \((1, 3)\)
\end{question}

%------------------------------------------------------------------------------%
\begin{resolution}{Solução}
  \onslide<+->{a) Determinando a curva de nível}
  \begin{align*}
    \onslide<+->{f(x, y)    & = 6}   \\
    \onslide<+->{6x -3y + 9 & = 6} \\
    \onslide<+->{-3y & = 6 -6x - 9} \\
    \onslide<+->{-3y & = -6x -3} \\
    \onslide<+->{  y & =  2x + 1}
  \end{align*}
  \onslide<+->{A curva de nível é uma reta}
\end{resolution}

%------------------------------------------------------------------------------%
\begin{resolution}{Solução}
  b) Calculando as derivadas parciais de $f$

  \pause
  \[
    \D{f}{x}
    \pause
    \;=\; \D{}{x}\left(6x - 3y + 9\right)
    \pause
    \;=\; 6
  \]

  \pause
  \[
    \D{f}{y}
    \pause
    \;=\; \D{}{y}\left(6x - 3y + 9\right)
    \pause
    \;=\; -3
  \]

  \pause
  Gradiente de $f$
  \pause
  \[
    \nabla f(x, y)
    \pause
    \;=\; \nabla\left(6x - 3y + 9\right)
    \pause
    \;=\; \vetor{6}{-3}
  \]
\end{resolution}

%------------------------------------------------------------------------------%
\begin{resolution}{Solução}
  Calculando o gradiente nos pontos
  \((0, 1)\) e
  \((1, 3)\)

  \pause
  \[
    \nabla f(0, 1)
    \pause
    \;=\; \vetor{6}{-3}
  \]

  \pause
  \[
    \nabla f(2, -1)
    \pause
    \;=\; \vetor{6}{-3}
  \]
\end{resolution}

%------------------------------------------------------------------------------%